\begin{frame}
  \onslide<1->{
  	\frametitle{Dynamic Systems}
	\framesubtitle{Measurement Models}  
 }
\begin{itemize}
  \onslide<2->{
	\item  In engineering applications, we often model not
	 only system dynamics but also the processes by which we \textbf{measure} system outputs
	}
	\begin{itemize}
		\onslide<3->{\item This is usually because we need to control the system dynamics
		}
		\onslide<4->{\item But more often than not we need to estimate certain quantities of interest associated with our system
		}
		\onslide<5->{\item And we can't control or change the behavior of something we can't accurately estimate!
		}
	\end{itemize}
\end{itemize}
\onslide<6->{\center --------------------------------------------------------------------}
\end{frame}

\begin{frame}
	\onslide<1->{
		\frametitle{Dynamic Systems}
		\framesubtitle{Measurement Models}  
	}
	\begin{itemize}
		\onslide<2->{
			\item  But outputs are not always directly observable!! They might have a complicated relationship to something we can physically measure
		}
		\onslide<3->{
			\item Thus, in modeling dynamic systems, we often add a second set of 
			equations, usually algebraic or functional in nature, that relates
			the output $y$ (and possibly the input $u$) to some measured or observed quantities
			\[z = h(y)+g(u)\]  
		}
	\end{itemize}
	\onslide<4->{\center --------------------------------------------------------------------}
\end{frame}

\begin{frame}
	\onslide<1->{\frametitle{State Space Representation Example} 
		\framesubtitle{Simple Harmonic Oscillator}    
	}
	\begin{itemize}
		\onslide<2->{
			\item  For a measurement model, we could consider a camera with a strobe flash that periodically records the instantaneous position of the sliding mass.
		}
		\onslide<3->{
			\item   In this case, only the position would be observable and we would write the measurement or observation model as:
			\[
			z =
			\begin{bmatrix}
			1 & 0
			\end{bmatrix}
			\begin{bmatrix}
			x_1\\
			x_2
			\end{bmatrix}.
			\]
			(Remember that we are ignoring both process and measurement noise.)   
		}
	\end{itemize}
	\onslide<4->{\center --------------------------------------------------------------------}
\end{frame}
